\documentclass[11pt,a4paper]{article}

% ============================ PACKAGES ============================
\usepackage[utf8]{inputenc}
\usepackage[T1]{fontenc}
\usepackage[english]{babel}
\usepackage{geometry}
\geometry{margin=2.5cm}
\usepackage{graphicx}
\usepackage{booktabs}
\usepackage{amsmath}
\usepackage{amssymb}
\usepackage{xcolor}
\usepackage{hyperref}
\hypersetup{
    colorlinks=true,
    linkcolor=blue!50!black,
    citecolor=blue!50!black,
    urlcolor=blue!60!black
}
\usepackage{caption}
\captionsetup{font=small,labelfont=bf}
\usepackage{booktabs}
\usepackage{array}
\usepackage{multirow}
\usepackage{enumitem}
\usepackage{fancyhdr}
\pagestyle{fancy}
\fancyhf{}
\rhead{\small UAV Segmentation Robustness to Synthetic Fog}
\lhead{\small}
\cfoot{\thepage}
\usepackage{abstract}
\usepackage{titlesec}
\titleformat{\section}{\normalfont\Large\bfseries}{\thesection}{1em}{}
\titleformat{\subsection}{\normalfont\large\bfseries}{\thesubsection}{1em}{}

% Placeholder figure command: draws a box if the image file is missing,
% so the document always compiles even before you add the PNGs.
\newcommand{\figordummy}[3]{%
  \IfFileExists{figures/#1}{%
    \includegraphics[width=#2]{figures/#1}%
  }{%
    \fbox{\parbox[c][4cm][c]{#2}{\centering{\ttfamily figures/#1}\par\vspace{4pt}{\small (#3)}}}%
  }%
}

% ============================ TITLE ============================
\title{\textbf{Evaluating and Improving the Robustness of UAV\\
Semantic Segmentation to GAN-Generated Fog}}
\author{Fabrizio Cozzolino\\
\small Computer Vision Project Report}
\date{\today}

\begin{document}

\maketitle
\thispagestyle{empty}

% ============================ ABSTRACT ============================
\begin{abstract}
\noindent
Semantic segmentation models for unmanned aerial vehicle (UAV) imagery are
typically trained on clear-weather data and degrade substantially under adverse
conditions such as fog. Because no public aerial dataset provides paired
clear/foggy images, we investigate a domain-transfer strategy: a Pix2Pix
conditional GAN is trained on the street-level Foggy Cityscapes dataset and
applied to the aerial Varied Drone Dataset (VDD) to synthesise foggy variants
at two severities. We first establish a U-Net baseline (test mIoU $0.7168$),
diagnose and fix a catastrophic class-imbalance failure on the \emph{vehicle}
class, then quantify the robustness drop under synthetic fog ($-7.2\%$ for
medium, $-25.0\%$ for dense). Retraining the U-Net on a balanced mix of clean
and foggy data recovers $82\%$ and $83\%$ of the lost performance respectively,
while \emph{improving} clean-data accuracy. Finally, a controlled comparison
against a purely algorithmic fog augmentation (Albumentations) reveals a
counterintuitive generalisation asymmetry: the GAN-augmented model excels on
GAN-style fog but collapses on algorithmic fog, whereas the
algorithmic-augmented model remains robust across both fog types and on clean
data. This suggests that, for robustness, a broader and more aggressive
perturbation distribution can outperform a realistic but narrow learned one.
All code, data and models are publicly released.
\end{abstract}

\vspace{0.5cm}
\noindent\textbf{Keywords:} semantic segmentation, UAV, domain adaptation,
generative adversarial networks, Pix2Pix, robustness, data augmentation, fog.

\newpage
\tableofcontents
\newpage

% ============================ 1. INTRODUCTION ============================
\section{Introduction}

\subsection{Motivation}
Unmanned aerial vehicles (UAVs) are increasingly deployed for tasks that
require dense scene understanding: infrastructure inspection, precision
agriculture, environmental monitoring, and search-and-rescue. Semantic
segmentation, which assigns a class label to every pixel, is a core enabling
technology for these applications. However, models trained on clear-weather
imagery tend to fail when deployed in adverse atmospheric conditions. Fog in
particular reduces contrast, attenuates distant structures, and shifts the
colour distribution of the scene, all of which violate the assumptions a
clear-trained network has learned.

For ground-level autonomous driving, this problem has been studied with
dedicated foggy datasets. For aerial imagery, no comparable public dataset of
paired clear/foggy scenes exists, which makes both the evaluation and the
mitigation of fog-induced degradation an open problem.

\subsection{Problem Statement}
We address two questions:
\begin{enumerate}[leftmargin=*]
  \item \textbf{Quantification:} how much does a clear-trained UAV
  segmentation model degrade under fog of increasing severity?
  \item \textbf{Mitigation:} can synthetic fog, generated by a model trained
  on a \emph{different} visual domain, be used as data augmentation to recover
  robustness?
\end{enumerate}

\subsection{Contributions}
This work makes the following contributions:
\begin{enumerate}[leftmargin=*]
  \item An end-to-end, reproducible PyTorch pipeline for UAV semantic
  segmentation under synthetic fog, spanning baseline training, GAN-based fog
  synthesis, robustness measurement and augmented retraining.
  \item A diagnosis and resolution of a class-imbalance failure that drove the
  \emph{vehicle} IoU to exactly zero, recovered to $0.585$ through a combined
  resolution and loss-weighting intervention.
  \item A quantitative characterisation of the robustness drop under
  GAN-generated fog ($-7.2\%$ at medium severity, $-25.0\%$ at dense).
  \item A demonstration that mixed clean+foggy retraining recovers $\approx
  82$--$83\%$ of the lost mIoU while slightly \emph{improving} clean
  performance.
  \item A controlled comparison between GAN-based and algorithmic fog
  augmentation that uncovers a counterintuitive generalisation asymmetry,
  with practical implications for robustness-oriented augmentation design.
\end{enumerate}

% ============================ 2. RELATED WORK ============================
\section{Related Work}

\paragraph{Semantic segmentation.}
Fully convolutional networks~\cite{long2015fcn} established dense per-pixel
prediction as an end-to-end learnable task. The U-Net
architecture~\cite{ronneberger2015unet} introduced a symmetric
encoder--decoder with skip connections that propagate high-resolution spatial
detail across the bottleneck, and remains a strong and widely adopted baseline,
particularly when training data is limited. Using an ImageNet-pretrained
backbone as the encoder (here ResNet-34~\cite{he2016resnet}) is standard
practice: transfer learning provides a strong initialisation and accelerates
convergence on small datasets. We use the
\texttt{segmentation\_models\_pytorch}~\cite{iakubovskii2019smp} implementation.

\paragraph{Aerial scene understanding.}
Drone imagery differs from ground-level imagery in viewpoint (near-nadir or
oblique), object scale (vehicles may span only tens of pixels), and class
distribution. The Varied Drone Dataset (VDD)~\cite{cai2023vdd} used here
provides pixel-level annotations across seven classes for high-resolution
aerial scenes.

\paragraph{Adverse-weather robustness.}
Foggy Cityscapes~\cite{sakaridis2018foggy} augments street-level
scenes~\cite{cordts2016cityscapes} with synthetic fog generated via
the Koschmieder optical model, which makes attenuation depth-dependent. Such
datasets have been used to study and improve robustness for autonomous driving;
the aerial counterpart is comparatively unexplored, motivating the
domain-transfer approach taken here.

\paragraph{Image-to-image translation.}
Conditional GANs~\cite{goodfellow2014gan,mirza2014cgan}, and
Pix2Pix~\cite{isola2017pix2pix} in particular, learn a mapping between paired
image domains using an adversarial loss combined with an $L_1$ reconstruction
term and a PatchGAN discriminator that judges local texture realism. We
repurpose Pix2Pix as a fog \emph{synthesiser}: it learns the clear$\rightarrow$
foggy transformation from paired data and is then applied out-of-domain.

\paragraph{Data augmentation for robustness.}
Augmentation is a standard defence against distribution shift. Our final study
contrasts a learned augmentation (the GAN) against a classical algorithmic one
(the parametric fog overlay of Albumentations~\cite{buslaev2020albumentations}),
and analyses which yields more transferable robustness.

% ============================ 3. DATASETS ============================
\section{Datasets}

\subsection{Varied Drone Dataset (VDD)}
VDD comprises 400 high-resolution aerial images ($4000\times3000$) annotated
into seven semantic classes: \emph{other, wall, road, vegetation, vehicle,
roof, water}. We use the authors' split of 280 training, 80 validation and 40
test images. The dataset exhibits severe class imbalance: as shown in
Table~\ref{tab:class_dist}, \emph{vehicle} occupies only $\approx 0.6\%$ of all
pixels despite appearing in $60.7\%$ of training images. This imbalance plays a
central role in Section~\ref{sec:vehicle}.

\begin{table}[ht]
\centering
\caption{Approximate pixel-frequency distribution of the seven VDD classes on
the training split. The \emph{vehicle} class is two orders of magnitude rarer
than \emph{vegetation}.}
\label{tab:class_dist}
\begin{tabular}{lc}
\toprule
\textbf{Class} & \textbf{Pixel frequency (\%)} \\
\midrule
vegetation & $\sim$30 \\
roof       & $\sim$20 \\
road       & $\sim$15 \\
water      & $\sim$10 \\
wall       & $\sim$8  \\
other      & $\sim$6  \\
vehicle    & $\sim$0.6 \\
\bottomrule
\end{tabular}
\end{table}

Figure~\ref{fig:vdd_samples} shows representative VDD scenes with their
ground-truth annotations. The near-nadir and oblique aerial viewpoints, the
small apparent size of vehicles, and the dominance of large homogeneous
regions (vegetation, roofs, water) are all visible and are properties that
distinguish this setting from ground-level segmentation.

\begin{figure}[ht]
\centering
\figordummy{vdd_samples.png}{0.92\textwidth}{VDD: image / colorised mask /
overlay, four examples}
\caption{Representative VDD samples: input aerial image (left), colourised
seven-class ground-truth mask (centre), and overlay (right). The dataset
spans dense urban blocks, water bodies, vegetation and mixed scenes.}
\label{fig:vdd_samples}
\end{figure}

\subsection{Foggy Cityscapes}
Foggy Cityscapes provides 500 street-level scenes, each available in three
paired versions: \emph{No\_Fog}, \emph{Medium\_Fog} and \emph{Dense\_Fog}. The
fog is synthetic, generated from the scene depth via the Koschmieder
scattering model, so attenuation increases with distance. The paired structure
(identical scene across fog levels) is exactly what Pix2Pix requires.

Figure~\ref{fig:foggy_cs} shows the two severities used in this work. The
third column, the normalised absolute difference $|\text{foggy}-\text{clean}|$,
makes the depth-dependence explicit: the effect is strongest on distant
structures and weakest in the near field, a physical property that the GAN
implicitly learns and that becomes relevant when the model is later applied to
near-nadir aerial views with little depth variation.

\begin{figure}[ht]
\centering
\begin{minipage}{0.49\textwidth}
\centering
\figordummy{foggy_cs_medium.png}{\textwidth}{Foggy Cityscapes medium}\par\vspace{2pt}
{\footnotesize (a) medium fog}
\end{minipage}\hfill
\begin{minipage}{0.49\textwidth}
\centering
\figordummy{foggy_cs_dense.png}{\textwidth}{Foggy Cityscapes dense}\par\vspace{2pt}
{\footnotesize (b) dense fog}
\end{minipage}
\caption{Foggy Cityscapes training pairs at the two severities. Columns:
clean (\emph{No\_Fog}), foggy target, and normalised difference. The
difference map reveals the depth-dependent nature of the synthetic fog
(stronger far, weaker near).}
\label{fig:foggy_cs}
\end{figure}

\subsection{The Domain-Transfer Rationale}
Since no aerial dataset offers paired clear/foggy images, we learn the fog
transformation where paired data \emph{does} exist (street level) and apply it
where it is needed (aerial). This deliberately introduces a domain gap between
the GAN's training distribution (Cityscapes) and its deployment distribution
(VDD); quantifying and discussing the consequences of this gap is part of the
study.

% ============================ 4. METHODOLOGY ============================
\section{Methodology}

The pipeline is organised in six phases. Phases 1--5 form the core study;
Phase~6 is a controlled extension.

\subsection{Phase 1 --- Baseline U-Net}
The segmentation model is an \texttt{smp.Unet} with a ResNet-34 encoder
pretrained on ImageNet (24.4M parameters), producing seven-channel logits at
the input resolution. Training uses AdamW~\cite{loshchilov2019adamw}
(learning rate $10^{-4}$, weight decay $10^{-4}$),
cosine-annealed~\cite{loshchilov2017sgdr} over 30 epochs, batch size 4,
gradient clipping at $1.0$, and a cross-entropy loss. Metrics are computed
with \texttt{torchmetrics}~\cite{torchmetrics}. Two design choices are critical
and motivated in Section~\ref{sec:vehicle}: an input resolution of
$768\times768$ and inverse-square-root class weighting in the loss.

\subsection{Phase 2 --- Pix2Pix Fog Generator}
The generator $G$ is a U-Net (ResNet-34 encoder) producing a three-channel RGB
image with a final $\tanh$ activation, so its output lies in $[-1,+1]$;
targets are normalised to the same range. The discriminator $D$ is a
$70\times70$ PatchGAN (5 strided convolutions, $\approx 2.77$M parameters,
$9\times$ smaller than $G$) that receives the channel-wise concatenation of
the conditioning (clear) image and the candidate foggy image, yielding a
$30\times30$ map of real/fake logits. The objective is
\begin{equation}
\mathcal{L}_{G} \;=\; \mathcal{L}_{\text{adv}} \;+\; \lambda_{L_1}\,
\mathcal{L}_{L_1}, \qquad \lambda_{L_1}=100,
\end{equation}
with $\mathcal{L}_{\text{adv}}$ a binary cross-entropy on the discriminator
logits and $\mathcal{L}_{L_1}=\lVert G(x)-y\rVert_1$. Training uses
Adam~\cite{kingma2015adam} with
$\beta=(0.5,0.999)$, learning rate $2\times10^{-4}$, at $256\times256$
resolution for 50 epochs. Two generators are trained independently, one per
fog severity (medium, dense).

\subsection{Phase 3 --- Foggy VDD Generation}
Each trained generator is applied to all 400 VDD images. Because the GAN was
trained at $256$ but the segmentation network operates at $768$, we generate at
both application scales, producing four datasets
(\{medium,dense\}$\times$\{256,768\}). All outputs are stored at
$768\times768$. A subtle but important requirement is that the segmentation
pipeline demands identical spatial dimensions for image and mask: masks are
therefore resized to $768\times768$ using nearest-neighbour interpolation,
which preserves the discrete class IDs (bilinear interpolation would create
invalid fractional labels).

\subsection{Phase 4 --- Robustness Measurement}
The Phase~1 model is evaluated, without any retraining, on the clean test set
and on each foggy variant. The drop in mIoU relative to the clean baseline
quantifies the model's fragility to fog.

\subsection{Phase 5 --- Mixed Retraining}
The U-Net is retrained from the same initialisation on a balanced
concatenation (\texttt{ConcatDataset}) of clean, foggy-medium and foggy-dense
training images (280 each, 840 total), validating on the clean split only so
the selection criterion is comparable to Phase~1. All other hyperparameters are
held fixed to isolate the effect of the data augmentation.

\subsection{Phase 6 --- Algorithmic Fog Baseline}
To assess whether the GAN is necessary, we generate fog directly on VDD with a
deterministic \texttt{Albumentations.RandomFog} transform at two intensities,
retrain a third model (\texttt{v3\_alg}) on the analogous clean+algorithmic mix,
and cross-evaluate all models on all fog types.

% ============================ 5. RESULTS ============================
\section{Experimental Results}

\subsection{Baseline and the Vehicle-Class Failure}
\label{sec:vehicle}
The first training run (\texttt{v1}: $512\times512$, no class weighting)
reached a validation mIoU of $0.6524$, an apparently reasonable figure.
Per-class inspection, however, revealed that the \emph{vehicle} IoU was
exactly $0.0000$: the model never predicted a single vehicle pixel.

Investigation showed a dual cause. First, \emph{vehicle} accounts for only
$\approx 0.6\%$ of pixels, so a cross-entropy loss is minimised almost as well
by ignoring the class entirely. Second, downsampling $4000\times3000$ images to
$512\times512$ shrinks a typical vehicle to a handful of pixels, effectively
erasing it. The fix combined two interventions: raising the input resolution to
$768\times768$ (reducing the area-reduction factor) and applying
inverse-square-root class weights in the cross-entropy (a $1/\sqrt{f}$ weighting
yields a moderate $\sim$3--5$\times$ emphasis on \emph{vehicle}, whereas pure
$1/f$ would over-weight it by $\sim$150$\times$ and destabilise training).

The effect is summarised in Table~\ref{tab:v1v2} and visible in the training
curves of Figures~\ref{fig:tb_p1_summary}--\ref{fig:tb_p1_vehicle}: the
\emph{vehicle} IoU rises from a flat zero to $0.585$, and overall validation
mIoU improves by more than ten points to $0.7602$. The official test mIoU of
the corrected model (\texttt{v2}) is $0.7168$; the $\approx 4$-point
validation--test gap is the expected consequence of model selection on the
validation set.

\begin{table}[ht]
\centering
\caption{Per-class IoU before (\texttt{v1}) and after (\texttt{v2}) the
class-imbalance fix (validation split). The \emph{vehicle} class is recovered
from exactly zero.}
\label{tab:v1v2}
\begin{tabular}{lcc}
\toprule
\textbf{Class} & \textbf{v1} & \textbf{v2} \\
\midrule
other        & 0.6217 & 0.6506 \\
wall         & 0.5214 & 0.5773 \\
road         & 0.6985 & 0.7347 \\
vegetation   & 0.8839 & 0.8898 \\
\textbf{vehicle} & \textbf{0.0000} & \textbf{0.5750} \\
roof         & 0.8899 & 0.9082 \\
water        & 0.9514 & 0.9583 \\
\midrule
\textbf{mIoU} & \textbf{0.6524} & \textbf{0.7602} \\
\bottomrule
\end{tabular}
\end{table}

\begin{figure}[ht]
\centering
\figordummy{tb_p1_summary.png}{0.9\textwidth}{Phase 1: val loss and val mIoU,
v1 (red) vs v2 (blue)}
\caption{Phase~1 validation curves. The corrected model \texttt{v2} (blue)
reaches a markedly higher mIoU and a lower loss plateau than the baseline
\texttt{v1} (red).}
\label{fig:tb_p1_summary}
\end{figure}

\begin{figure}[ht]
\centering
\figordummy{tb_p1_vehicle.png}{0.9\textwidth}{val iou\_vehicle / wall / water:
v1 flat at zero, v2 rising}
\caption{The \emph{vehicle} IoU (left panel). The baseline \texttt{v1} (red)
remains flat at zero for the entire training, while \texttt{v2} (blue) rises to
$\approx 0.58$, confirming that the combined resolution and loss-weighting fix
resolves the failure.}
\label{fig:tb_p1_vehicle}
\end{figure}

\subsection{Pix2Pix Training}
Both generators were trained for 50 epochs ($\approx 57$ minutes each on a
Tesla T4). The best validation $L_1$ was $0.0385$ (medium) and $0.0468$
(dense), both at epoch 38. The $L_1$ reconstruction term decreased
monotonically and stably (Figure~\ref{fig:tb_p2_l1}), while the adversarial
dynamics oscillated: the discriminator periodically ``won'' (its loss
collapsing toward zero with $D(\text{real})\rightarrow1$,
$D(\text{fake})\rightarrow0$) before the equilibrium was restored
(Figure~\ref{fig:tb_p2_dyn}). Because the best checkpoint is selected on
validation $L_1$ (Figure~\ref{fig:tb_p2_vall1}), the saved generator
corresponds to a balanced epoch rather than a late one in which $D$ dominated;
the heavy $\lambda_{L_1}=100$ weighting also keeps generation faithful even
when the adversarial signal becomes uninformative.

\begin{figure}[ht]
\centering
\figordummy{tb_p2_dyn.png}{0.95\textwidth}{Pix2Pix D\_fake\_mean, D\_real\_mean,
loss\_D --- oscillating adversarial dynamics}
\caption{Adversarial dynamics of the two Pix2Pix runs (medium red, dense blue).
The oscillation between equilibrium and discriminator dominance is the expected
behaviour of an unregularised GAN game.}
\label{fig:tb_p2_dyn}
\end{figure}

\begin{figure}[ht]
\centering
\begin{minipage}{0.48\textwidth}
\centering
\figordummy{tb_p2_g_l1.png}{\textwidth}{loss\_G\_l1 monotone decrease}
\end{minipage}\hfill
\begin{minipage}{0.48\textwidth}
\centering
\figordummy{tb_p2_val_l1.png}{\textwidth}{val/L1, best at epoch 38}
\end{minipage}
\caption{Left: the generator $L_1$ term decreases smoothly despite the
adversarial oscillation. Right: validation $L_1$, the model-selection
criterion; the best checkpoint is epoch 38 for both severities.}
\label{fig:tb_p2_l1}
\hphantom{}
\label{fig:tb_p2_vall1}
\end{figure}

The qualitative effect of this training is shown in
Figure~\ref{fig:pix2pix_progression}. At epoch~1 the generator produces
incoherent, artefact-ridden output (saturated colour blobs, heavy blur); by
epoch~50 the synthesised foggy images (middle row) closely match the
ground-truth targets (bottom row), reproducing both the photometric haze and
the depth-dependent attenuation of the reference fog. The dense generator
faces a harder transformation and its epoch-1 output is correspondingly more
degraded, but it too converges to a convincing result.

\begin{figure}[ht]
\centering
\begin{minipage}{0.48\textwidth}
\centering
\figordummy{pix2pix_med_e001.png}{\textwidth}{medium, epoch 1}\par\vspace{2pt}
{\footnotesize (a) medium fog --- epoch 1}
\end{minipage}\hfill
\begin{minipage}{0.48\textwidth}
\centering
\figordummy{pix2pix_med_e050.png}{\textwidth}{medium, epoch 50}\par\vspace{2pt}
{\footnotesize (b) medium fog --- epoch 50}
\end{minipage}

\vspace{6pt}

\begin{minipage}{0.48\textwidth}
\centering
\figordummy{pix2pix_dns_e001.png}{\textwidth}{dense, epoch 1}\par\vspace{2pt}
{\footnotesize (c) dense fog --- epoch 1}
\end{minipage}\hfill
\begin{minipage}{0.48\textwidth}
\centering
\figordummy{pix2pix_dns_e050.png}{\textwidth}{dense, epoch 50}\par\vspace{2pt}
{\footnotesize (d) dense fog --- epoch 50}
\end{minipage}
\caption{Pix2Pix learning progression. In each grid the rows are
\emph{clean} input (top), \emph{generated} foggy (middle) and \emph{real}
foggy target (bottom). The generator evolves from incoherent output at
epoch~1 (a,c) to a faithful reproduction of the target fog at epoch~50 (b,d),
for both medium and dense severities.}
\label{fig:pix2pix_progression}
\end{figure}

\subsection{Generated Foggy VDD --- Qualitative Assessment}
Figure~\ref{fig:vdd_comp} shows the four GAN-generated variants alongside the
clean image. The generator successfully transfers the photometric appearance of
fog --- reduced contrast, whitened background, softened distant structures ---
while preserving object identities. A limitation is visible: because the
training fog (Foggy Cityscapes) is depth-dependent, but near-nadir aerial
views have little depth variation, the applied fog is comparatively uniform on
aerial scenes, which is physically less accurate than ground-level fog. The
$256$-application variant additionally suffers from upscaling artefacts, which
the quantitative results in Section~\ref{sec:robustness} show to be dominant.

\begin{figure}[ht]
\centering
\figordummy{vdd_foggy_comp.png}{\textwidth}{VDD clean vs 4 GAN foggy variants}
\caption{Clean VDD test images (column~1) and the four GAN-generated foggy
variants. Fog severity and the medium/dense distinction are visible; the
near-uniform application on aerial views (no strong depth gradient) is the
expected consequence of the street-level$\rightarrow$aerial domain gap.}
\label{fig:vdd_comp}
\end{figure}

\subsection{Robustness Drop}
\label{sec:robustness}
Table~\ref{tab:phase4} reports the Phase~1 model evaluated on all variants.
Two distinct stories emerge. The $256$-application variants show catastrophic
drops ($-63.7\%$, $-70.9\%$); however, Figure~\ref{fig:phase4} and visual
inspection attribute this to upscaling artefacts that dominate the fog effect,
not to fog itself. We therefore treat the $256$ results as an ablation that
\emph{justifies} working at $768$, and report the $768$ variants as the
official measurement: a modest $-7.2\%$ under medium fog and a substantial
$-25.0\%$ under dense fog. Per-class analysis shows that large, homogeneous
surfaces (\emph{roof}, \emph{road}, \emph{water}) degrade most, while textured
classes (\emph{vegetation}, \emph{vehicle}) are comparatively resilient.

\begin{table}[ht]
\centering
\caption{Phase~4 robustness of the clean-trained U-Net (\texttt{v2}) on the
test set. The @256 variants are invalidated by upscaling artefacts and
reported only as an ablation.}
\label{tab:phase4}
\begin{tabular}{lcc}
\toprule
\textbf{Test set} & \textbf{mIoU} & \textbf{$\Delta$ vs clean} \\
\midrule
clean              & 0.7168 & --- \\
GAN medium @256    & 0.2601 & $-63.7\%$ \emph{(ablation)} \\
\textbf{GAN medium @768} & \textbf{0.6652} & $\mathbf{-7.2\%}$ \\
GAN dense @256     & 0.2086 & $-70.9\%$ \emph{(ablation)} \\
\textbf{GAN dense @768}  & \textbf{0.5377} & $\mathbf{-25.0\%}$ \\
\bottomrule
\end{tabular}
\end{table}

\begin{figure}[ht]
\centering
\figordummy{phase4_per_class.png}{0.95\textwidth}{Per-class IoU, clean vs 4
foggy variants}
\caption{Per-class IoU for the clean-trained model across all variants. The
@256 bars (light blue, orange) are uniformly depressed by upscaling artefacts;
the @768 bars (blue, dark red) reflect the genuine fog effect, with large
homogeneous classes most affected.}
\label{fig:phase4}
\end{figure}

\subsection{Recovery via Mixed Retraining}
Table~\ref{tab:phase5} and Figure~\ref{fig:phase5} report the central result.
Retraining on the clean+foggy mix (\texttt{v3}) recovers $82\%$ of the medium
loss and $83\%$ of the dense loss, and --- crucially --- does \emph{not}
sacrifice clean performance: clean mIoU rises slightly from $0.7168$ to
$0.7264$. The augmentation therefore behaves as a beneficial regulariser
rather than a specialisation trade-off. The validation curves
(Figure~\ref{fig:tb_p5_summary}) confirm a higher mIoU plateau and a lower loss
for \texttt{v3}.

\begin{table}[ht]
\centering
\caption{Phase~5: clean-trained \texttt{v2} vs mixed-trained \texttt{v3}
(test mIoU). Recovery is computed relative to the clean baseline.}
\label{tab:phase5}
\begin{tabular}{lcccc}
\toprule
\textbf{Test set} & \textbf{v2} & \textbf{v3} & \textbf{$\Delta$} &
\textbf{Recovery} \\
\midrule
clean        & 0.7168 & \textbf{0.7264} & $+0.0097$ & --- \\
GAN medium   & 0.6652 & \textbf{0.7076} & $+0.0424$ & $\sim$82\% \\
GAN dense    & 0.5377 & \textbf{0.6866} & $+0.1490$ & $\sim$83\% \\
\bottomrule
\end{tabular}
\end{table}

\begin{figure}[ht]
\centering
\figordummy{phase5_v2_v3.png}{0.75\textwidth}{Bar chart v2 vs v3 on
clean/medium/dense}
\caption{Robustness recovery via foggy data augmentation. The mixed-trained
\texttt{v3} (green) outperforms the clean-trained \texttt{v2} (blue) on every
test set, with the largest gain on the hardest (dense) condition.}
\label{fig:phase5}
\end{figure}

\begin{figure}[ht]
\centering
\figordummy{tb_p5_summary.png}{0.9\textwidth}{Phase 5 val loss and val mIoU,
v2 vs v3}
\caption{Phase~5 validation curves (clean validation split). \texttt{v3}
(blue) attains a higher mIoU plateau and a lower loss than \texttt{v2} (red),
despite being trained on a harder, more varied mixture.}
\label{fig:tb_p5_summary}
\end{figure}

Per-class recovery on the hardest condition (GAN dense) is shown in
Table~\ref{tab:phase5_perclass}: the classes that suffered most in Phase~4
(\emph{roof}, \emph{water}) are precisely those that recover most, indicating
that the augmentation repairs the specific weaknesses of the clean-trained
model.

\begin{table}[ht]
\centering
\caption{Per-class IoU on GAN dense @768: clean-trained \texttt{v2} vs
mixed-trained \texttt{v3}.}
\label{tab:phase5_perclass}
\begin{tabular}{lcc}
\toprule
\textbf{Class} & \textbf{v2} & \textbf{v3} \\
\midrule
other       & 0.3568 & \textbf{0.5348} \\
wall        & 0.5233 & \textbf{0.5873} \\
road        & 0.4650 & \textbf{0.6067} \\
vegetation  & 0.7647 & \textbf{0.8596} \\
vehicle     & 0.4789 & \textbf{0.4993} \\
roof        & 0.5249 & \textbf{0.8038} \\
water       & 0.6501 & \textbf{0.9148} \\
\midrule
\textbf{mIoU} & 0.5377 & \textbf{0.6866} \\
\bottomrule
\end{tabular}
\end{table}

% ============================ 6. EXTENSION ============================
\section{Extension: GAN vs Algorithmic Fog}

\subsection{Motivation and Setup}
Given the domain gap incurred by training the GAN on street-level data, a
natural question is whether a simple algorithmic fog applied \emph{directly}
to VDD would serve as well, or better. We generate medium and dense fog with a
deterministic \texttt{Albumentations.RandomFog} transform and retrain a third
model (\texttt{v3\_alg}) on the analogous clean+algorithmic mix, then
cross-evaluate all three models on all fog types.

Figure~\ref{fig:gan_alg} contrasts the two fog sources visually. The
algorithmic fog is markedly different from the GAN fog: it combines a patchy
haze overlay with a pronounced defocus, making it visually more aggressive and
less photo-realistic than the GAN's smooth attenuation.

\begin{figure}[ht]
\centering
\figordummy{phase6_gan_alg.png}{\textwidth}{clean / GAN-med / alg-med /
GAN-dense / alg-dense}
\caption{Clean VDD vs GAN-generated and algorithmic fog at two severities. The
algorithmic fog (columns 3 and 5) is patchier and more defocused than the
smoother GAN fog (columns 2 and 4).}
\label{fig:gan_alg}
\end{figure}

\subsection{Results and the Generalisation Asymmetry}
Table~\ref{tab:phase6} reports the full cross-evaluation. The result is
counterintuitive. The GAN-augmented model \texttt{v3\_gan} is the best on every
GAN-fog test set, but \emph{collapses} on algorithmic fog (mIoU $0.117$ /
$0.106$), barely above the clean-trained baseline. By contrast, the
algorithmic-augmented model \texttt{v3\_alg} not only excels on algorithmic fog
(its training distribution) but also retains substantial performance on
GAN fog ($0.662$ / $0.607$) \emph{and} matches the baseline on clean data
($0.717$). In other words, \texttt{v3\_alg} generalises across both fog types,
whereas \texttt{v3\_gan} is specialised to its own.

\begin{table}[ht]
\centering
\caption{Full cross-evaluation (test mIoU). Bold marks the best model per row.
The clean-trained \texttt{v2} and the GAN-trained \texttt{v3\_gan} both
collapse on algorithmic fog, while \texttt{v3\_alg} stays robust across all
conditions.}
\label{tab:phase6}
\begin{tabular}{lccc}
\toprule
\textbf{Test set} & \textbf{v2 (clean)} & \textbf{v3\_gan} &
\textbf{v3\_alg} \\
\midrule
clean            & 0.7168 & \textbf{0.7264} & 0.7169 \\
GAN fog medium   & 0.6652 & \textbf{0.7076} & 0.6619 \\
GAN fog dense    & 0.5377 & \textbf{0.6866} & 0.6066 \\
Alg.\ fog medium & 0.0830 & 0.1174          & \textbf{0.6698} \\
Alg.\ fog dense  & 0.0708 & 0.1055          & \textbf{0.6652} \\
\bottomrule
\end{tabular}
\end{table}

% ============================ 7. DISCUSSION ============================
\section{Discussion}

\paragraph{Why does algorithmic fog generalise better?}
The GAN learns a narrow, realistic fog distribution: smooth attenuation with a
characteristic colour shift. A model trained on it becomes excellent at
inverting \emph{that specific} corruption but has never seen the heavy defocus
and patchiness of the algorithmic fog, so it fails out-of-distribution. The
algorithmic fog, conversely, is a broader and more aggressive perturbation
(haze \emph{and} blur \emph{and} spatial variability); training on it forces
the network to rely on more corruption-invariant features, which transfer even
to the smoother GAN fog. This is a concrete instance of the general principle
that, for robustness, the \emph{breadth} of the augmentation distribution can
matter more than its realism.

\paragraph{Domain gap.}
The street-level$\rightarrow$aerial transfer is imperfect: the depth-dependent
Koschmieder fog learned on Cityscapes becomes near-uniform on near-nadir aerial
scenes. This does not invalidate the GAN approach --- it still produces a
useful, recoverable corruption --- but it does limit the physical realism of
the synthesised aerial fog, and likely contributes to the narrowness of the
learned distribution.

\paragraph{Practical implication.}
If the operational goal is robustness to \emph{unknown} fog, a simple,
aggressive algorithmic augmentation is a strong and cheap baseline that may
outperform a more expensive, more realistic GAN. The GAN remains preferable if
the deployment fog is known to resemble the training fog.

\paragraph{Limitations.}
The study uses a single aerial dataset and a single backbone. No depth map is
available for VDD, so physically-grounded aerial fog could not be synthesised.
Hyperparameters were fixed rather than searched, and the Pix2Pix models were
trained for 50 epochs (vs the 100--200 typical in the literature) under compute
constraints. Statistical significance over the 40-image test set was not
formally established. These are natural avenues for future work.

% ============================ 8. CONCLUSIONS ============================
\section{Conclusions}

We presented an end-to-end pipeline that quantifies and mitigates the
fog-induced degradation of a UAV semantic segmentation model. After
diagnosing and resolving a class-imbalance failure (vehicle IoU $0\rightarrow
0.585$), we measured a robustness drop of $-7.2\%$ (medium) and $-25.0\%$
(dense) under GAN-generated fog, and showed that mixed clean+foggy retraining
recovers $\approx 82$--$83\%$ of the loss while slightly improving clean
accuracy. A controlled comparison against algorithmic fog augmentation revealed
a counterintuitive generalisation asymmetry: the realistic but narrow GAN fog
yields a specialised model, whereas the cruder but broader algorithmic fog
yields a more transferable one. For robustness-oriented augmentation, breadth
can beat realism.

Future work includes depth-aware aerial fog synthesis, validation across
multiple datasets and backbones, a hyperparameter grid search, longer GAN
training, and a formal statistical analysis of the observed differences.

% ============================ APPENDICES ============================
\appendix
\newpage

\section{Hyperparameter Configurations}
\begin{table}[ht]
\centering
\caption{Full training configurations.}
\begin{tabular}{lll}
\toprule
\textbf{Parameter} & \textbf{U-Net (v2/v3)} & \textbf{Pix2Pix} \\
\midrule
Backbone            & ResNet-34 (ImageNet) & ResNet-34 (ImageNet) \\
Input resolution    & $768\times768$       & $256\times256$ \\
Optimiser           & AdamW                & Adam $\beta=(0.5,0.999)$ \\
Learning rate       & $10^{-4}$            & $2\times10^{-4}$ \\
Weight decay        & $10^{-4}$            & --- \\
LR schedule         & Cosine annealing     & Constant \\
Epochs              & 30                   & 50 \\
Batch size          & 4                    & 8 \\
Gradient clipping   & 1.0                  & --- \\
Loss                & Weighted CE (inv.\ sqrt) & BCE $+\,100\,L_1$ \\
\bottomrule
\end{tabular}
\end{table}

\section{Reproducibility}
All code is publicly available at
\url{https://github.com/FabrizioCozzolino/fog-uav-robustness}. The pipeline was
developed locally and executed on Google Colab (Tesla T4). Datasets, trained
checkpoints and result files are versioned on the project storage. Each
training run records its configuration, per-epoch history and TensorBoard logs.

\section{Lessons Learned (Bug Log)}
Two non-trivial issues shaped the project and are documented here as part of
the scientific process:
\begin{enumerate}[leftmargin=*]
  \item \textbf{Vehicle IoU $=0$.} A $0.6\%$-frequency class combined with an
  aggressive $512\times512$ resize made the model ignore vehicles entirely.
  Resolved with $768\times768$ input and inverse-sqrt class weighting.
  \item \textbf{Mask shape mismatch.} Foggy images were saved at
  $768\times768$ but masks were initially copied at the original
  $4000\times3000$, violating the equal-dimension requirement of the
  augmentation pipeline. Resolved by nearest-neighbour mask resizing, chosen
  over bilinear to preserve discrete class IDs.
\end{enumerate}

\section{Complete Training Curves}
This appendix collects the full TensorBoard curves for completeness.

\paragraph{D.1 --- Phase 1 (U-Net baseline vs v2).}
\begin{figure}[ht]
\centering
\figordummy{tb_p1_train.png}{0.9\textwidth}{P1 train loss/lr}
\figordummy{tb_p1_val_a.png}{0.9\textwidth}{P1 val F1/acc/iou\_other}
\figordummy{tb_p1_val_b.png}{0.9\textwidth}{P1 val iou road/roof/vegetation}
\caption{Phase~1 training and validation curves (\texttt{v1} red,
\texttt{v2} blue).}
\end{figure}

\paragraph{D.2 --- Phase 2 (Pix2Pix medium vs dense).}
\begin{figure}[ht]
\centering
\figordummy{tb_p2_losses.png}{0.9\textwidth}{P2 D\_fake/D\_real/G\_adv}
\figordummy{tb_p2_batch.png}{0.9\textwidth}{P2 per-batch losses}
\caption{Phase~2 adversarial training curves (medium red, dense blue).}
\end{figure}

\paragraph{D.3 --- Phase 5 (v2 vs v3 mixed).}
\begin{figure}[ht]
\centering
\figordummy{tb_p5_train.png}{0.9\textwidth}{P5 train loss/lr}
\figordummy{tb_p5_val_a.png}{0.9\textwidth}{P5 val F1/acc/iou\_other}
\figordummy{tb_p5_val_b.png}{0.9\textwidth}{P5 val iou road/roof/vegetation}
\figordummy{tb_p5_val_c.png}{0.9\textwidth}{P5 val iou vehicle/wall/water}
\caption{Phase~5 training and validation curves (\texttt{v2} red,
\texttt{v3} blue).}
\end{figure}

% ============================ BIBLIOGRAPHY ============================
\newpage
\begin{thebibliography}{99}
\setlength{\itemsep}{2pt}

\bibitem{long2015fcn}
J.~Long, E.~Shelhamer, and T.~Darrell,
``Fully Convolutional Networks for Semantic Segmentation,''
in \emph{IEEE Conf. on Computer Vision and Pattern Recognition (CVPR)}, 2015.

\bibitem{ronneberger2015unet}
O.~Ronneberger, P.~Fischer, and T.~Brox,
``U-Net: Convolutional Networks for Biomedical Image Segmentation,''
in \emph{Medical Image Computing and Computer-Assisted Intervention
(MICCAI)}, 2015.

\bibitem{he2016resnet}
K.~He, X.~Zhang, S.~Ren, and J.~Sun,
``Deep Residual Learning for Image Recognition,''
in \emph{IEEE Conf. on Computer Vision and Pattern Recognition (CVPR)}, 2016.

\bibitem{iakubovskii2019smp}
P.~Iakubovskii,
``Segmentation Models PyTorch,''
\url{https://github.com/qubvel/segmentation_models.pytorch}, 2019.

\bibitem{cai2023vdd}
W.~Cai \emph{et al.},
``VDD: Varied Drone Dataset for Semantic Segmentation,''
arXiv preprint, 2023.

\bibitem{cordts2016cityscapes}
M.~Cordts \emph{et al.},
``The Cityscapes Dataset for Semantic Urban Scene Understanding,''
in \emph{IEEE Conf. on Computer Vision and Pattern Recognition (CVPR)}, 2016.

\bibitem{sakaridis2018foggy}
C.~Sakaridis, D.~Dai, and L.~Van~Gool,
``Semantic Foggy Scene Understanding with Synthetic Data,''
\emph{International Journal of Computer Vision}, vol.~126, 2018.

\bibitem{goodfellow2014gan}
I.~Goodfellow \emph{et al.},
``Generative Adversarial Nets,''
in \emph{Advances in Neural Information Processing Systems (NeurIPS)}, 2014.

\bibitem{mirza2014cgan}
M.~Mirza and S.~Osindero,
``Conditional Generative Adversarial Nets,''
arXiv preprint arXiv:1411.1784, 2014.

\bibitem{isola2017pix2pix}
P.~Isola, J.-Y.~Zhu, T.~Zhou, and A.~A.~Efros,
``Image-to-Image Translation with Conditional Adversarial Networks,''
in \emph{IEEE Conf. on Computer Vision and Pattern Recognition (CVPR)}, 2017.

\bibitem{kingma2015adam}
D.~P.~Kingma and J.~Ba,
``Adam: A Method for Stochastic Optimization,''
in \emph{International Conference on Learning Representations (ICLR)}, 2015.

\bibitem{loshchilov2019adamw}
I.~Loshchilov and F.~Hutter,
``Decoupled Weight Decay Regularization,''
in \emph{International Conference on Learning Representations (ICLR)}, 2019.

\bibitem{loshchilov2017sgdr}
I.~Loshchilov and F.~Hutter,
``SGDR: Stochastic Gradient Descent with Warm Restarts,''
in \emph{International Conference on Learning Representations (ICLR)}, 2017.

\bibitem{buslaev2020albumentations}
A.~Buslaev \emph{et al.},
``Albumentations: Fast and Flexible Image Augmentations,''
\emph{Information}, vol.~11, no.~2, 2020.

\bibitem{torchmetrics}
N.~Detlefsen \emph{et al.},
``TorchMetrics --- Measuring Reproducibility in PyTorch,''
\emph{Journal of Open Source Software}, 2022.

\end{thebibliography}

\end{document}
